\section{Methods}
\label{sec:methods}

\subsection{Architecture}

The codebase follows a layered structure (see \texttt{docs/architecture.md}):
\texttt{src/data} for dataset loading, registry, and leakage-safe preprocessing;
\texttt{src/quantum} for variational circuits (PennyLane) and hybrid models;
\texttt{src/classical} for matched baselines; and
\texttt{src/training} for the universal trainer, statistics, holdout evaluation, and logging.
Experiment orchestration stays thin in \texttt{experiments/exp\_NNN/run.py};
hyperparameters live in \texttt{config/experiments.yaml}.

Key data modules enforce split-before-scale ordering:
\texttt{splits.py} (stratified 70/30), \texttt{scaling.py} (train-only fit),
\texttt{real\_datasets.py} (UCI and MNIST PCA loaders), and \texttt{poisoning.py}
for robustness stress tests.

\subsection{Model contract}

All models implement \texttt{train}, \texttt{predict}, and \texttt{evaluate}.
Metrics are logged exclusively through \texttt{ExperimentLogger} to
\texttt{logs/experiments.jsonl} (append-only, never deleted).

\subsection{Evaluation protocol}

\begin{enumerate}
  \item Stratified 70/30 train--test split \emph{before} scaling or augmentation.
  \item Training on train set only; holdout accuracy on test set.
  \item Ten random seeds per publication profile; bootstrap 95\% CIs on mean accuracy.
  \item Paired Wilcoxon signed-rank tests with Holm-Bonferroni correction and Cohen's~$d$
        effect sizes where pairwise comparisons are reported.
\end{enumerate}

\subsection{Hyperparameters}

Default hyperparameters live in \texttt{config/experiments.yaml}.
Profiles: \texttt{publication} ($n{=}500$, 10 seeds), \texttt{publication\_large} ($n{=}1000$),
and \texttt{ci} (fast smoke, 1 seed).
Barren-plateau diagnostics (exp\_006) and adaptive learning-rate schedules (exp\_015)
use gradient-variance telemetry from \texttt{src/training/adaptive\_lr.py}.

\subsection{Conventional tabular comparison (exp\_058)}

Experiment \texttt{exp\_058} isolates whether the shipped \textbf{LargeNanoMLP}
(\texttt{exp\_032}) beats standard non-quantum stacks on open HIGGS tabular data.
Both deep models share the \emph{same parameter budget} ($\sim$1.14\,M weights,
topology $28{\rightarrow}2048{\rightarrow}512{\rightarrow}64{\rightarrow}1$); the
comparison is therefore a \textbf{training-stack ablation}, not a capacity mismatch.

\paragraph{Data protocol.}
We use the preprocessed \texttt{higgs\_v1} parquet splits (805K train / 172.5K val /
172.5K test). A \texttt{StandardScaler} is fit on train only and applied to val.
Seed~42 controls any row subsampling in CI profiles. Test rows are never used for
model selection.

\paragraph{LargeNanoMLP (our model).}
Implemented in PyTorch (\texttt{src/classical/large\_nano\_mlp.py}), trained in
\texttt{exp\_032} with batched Adam ($\mathrm{lr}{=}10^{-3}$, batch 2048, 12 epochs,
weight decay $10^{-4}$, dropout 0.3, binary cross-entropy). For \texttt{exp\_058} we
load the frozen shipped checkpoint (\texttt{dist/serve\_models/large\_nano\_mlp\_higgs/best.pt})
and report validation ROC-AUC without retraining.

\paragraph{Conventional baselines.}
We retrain four sklearn/XGBoost models on the \emph{same} scaled train matrix:
(i)~\texttt{LogisticRegression} (linear baseline);
(ii)~\texttt{MLPClassifier} with \texttt{hidden\_layer\_sizes=(2048,512,64)} --- the
primary matched deep competitor;
(iii)~\texttt{HistGradientBoostingClassifier} (default GBDT on CPU);
(iv)~shallow \texttt{XGBClassifier} (depth~3, 50 trees).
All baselines use seed~42; hyperparameters are documented in
\texttt{config/experiments.yaml} under \texttt{exp\_058\_conventional\_higgs\_baselines}.

\paragraph{Success gate.}
Pre-registered in \texttt{hypothesis.md}: LargeNanoMLP val ROC-AUC $\geq$ best
conventional + \textbf{0.5 percentage points}. Single-seed infrastructure comparison;
multi-seed Wilcoxon deferred.

\subsection{Reproducibility artifacts}

Figures and tables are regenerated from JSONL logs via \texttt{make figures} and
\texttt{make latex-tables}. Release bundles include SHA-256 checksums (\texttt{MANIFEST.txt})
for Zenodo archival (see \texttt{docs/zenodo.md}).
